\section*{Как оценивались вершина и две ветви}
При фиксированных \(w,T\) относительные свободные энергии восстанавливаются из гистограмм числа атомов Al \(k\) и химического порядка
\[
Q=\frac{4(4\max_j n_j-k)}{3N}.
\]
Величина \(n_j\) задаёт заселённость алюминием одной из четырёх подрешёток ГПУ. При данном химическом поле \(\Delta\mu\) и разделителе \(q_c\) сравниваются бассейны \(Q>q_c\) и \(Q\le q_c\). Четыре трансляционных варианта D0$_{19}$ дают одинаковый вклад. Для одной упорядоченной фазы относительно \(\alpha\) логарифм веса равен
\[
 D(T,\Delta\mu;w)=\ln\frac{P(Q>q_c)}{4P(Q\le q_c)},\qquad
 S(T;w)=\max_{\Delta\mu}D(T,\Delta\mu;w).
\]
На изотерме контакты ищутся по пересечениям \(D=0\) внутри прямой сетки полей. Для каждого поля контакта составы получаются как условные средние \(k/N\). Если оба края прямой сетки остаются внутри упорядоченной области, вычислить ширины обеих ветвей по этой сетке нельзя. Положение внутреннего максимума проверяется отдельно. В нём \(\partial_{\Delta\mu}D=\beta N(x_{\alpha_2}-x_\alpha)=0\). Конгруэнтной вершине конечной ячейки соответствует касание \(S=0\) при сохранении двух раздельных структурных мод.

При одинаковых \(T,\Delta\mu\) изменение разности полубольших потенциалов на атом при сужении стенки равно
\[
 \delta(\omega_{\alpha_2}-\omega_{\alpha})=
 -\frac{k_{\rm B}T}{N}[D(T,\Delta\mu;0{,}35)-D(T,\Delta\mu;0{,}40)].
\]
Величина относится к двум ограниченным по \(Q\) бассейнам конечной ячейки. Отрицательный знак означает увеличение относительного веса D0$_{19}$ при сужении стенки. При глубоком заходе в область устойчивости одной фазы барьер от слабой метастабильной моды может стать малым. Для сравнения потенциалов взяты общие непосредственно приложенные поля, для составов и теплоты --- собственные равновесные контакты и вершины каждого ансамбля.

Веса зонтиков и нормировки температур оцениваются методом MBAR внутри каждого \(w\) по исходным энергиям, \(k\) и \(Q\). Обычно выбран \(q_c=0{,}37\); проверяются и другие разделители между пиками. Изменение \(T_*\) вычисляется по разности двух независимо полученных оценок для \(w=0{,}35\) и \(0{,}40\). Статистические интервалы условны по ячейке, разделителю и покрытию диапазонов температур и поля. Для оценки доли повторов с двумя контактами учитываются и те блочные выборки, в которых границы не разрешились. Интервалы составов и ширин вычисляются по повторам с соответствующим корнем. \href{https://doi.org/10.1063/1.2978177}{Shirts и Chodera (2008).}
